\documentclass[10pt,a4paper]{amsart}
\usepackage[margin=19mm]{geometry}
\usepackage{amsmath,amssymb,mathtools}
\usepackage[scr=rsfs,cal=euler]{mathalfa}
\usepackage{xfrac}
\usepackage[unicode,hidelinks,bookmarks=false]{hyperref}
\newcommand{\ie}{i.e.\ }
\newcommand{\cf}{cf.\ }
\newcommand{\resp}{respectively\ }
\newcommand{\Subsection}{\subsection}
\providecommand{\nobreakdash}{\nobreak\hbox{-}}
\setlength{\emergencystretch}{2em}
\multlinegap0pt
\usepackage{bm}
\usepackage{bbm}
\usepackage{dsfont}
\usepackage{xspace}
\usepackage{cancel}
\usepackage{mathtools}
\usepackage{amsthm}
\makeatletter
\@ifundefined{theorem}{\newtheorem{theorem}{Theorem}}{}
\@ifundefined{lemma}{\newtheorem{lemma}[theorem]{Lemma}}{}
\makeatother
\renewcommand{\Im}{\operatorname{Im}}
\renewcommand{\Re}{\operatorname{Re}}
\title[Genus three embedded doubly periodic minimal surfaces with p]{Genus three embedded doubly periodic minimal surfaces with parallel ends}
\author{Peter Connor, Shoichi Fujimori, Phillip Marmorino,, Toshihiro Shoda}
\date{Source: arXiv:2105.10711v3 (CC BY 4.0)}
\begin{document}
\maketitle

\noindent\emph{Source and licence.} Genus three embedded doubly periodic minimal surfaces with parallel ends. Version arXiv:2105.10711v3 (CC BY 4.0), distributed under the Creative Commons Attribution 4.0 International licence, as recorded in the arXiv licence metadata for this exact version and shown on the abstract page https://arxiv.org/abs/2105.10711v3. Full rights notice: licenses/arxiv-2105.10711-CC-BY-4.0.txt.\par\smallskip

\noindent\textit{Editorial scope.} Counted unit: the Lemma whose source label is \texttt{lm:ReImw} in the article (source0.tex lines 642--649, source label \texttt{lm:ReImw}), delivered together with the article's own complete original proof of it (source0.tex lines 651--700, one \texttt{proof} environment of 50 lines). No mathematics is written, repaired or re-derived: statement and proof are the article's own text line for line. Required context retained uncounted: none. Every definition, display and label of those lines is retained unchanged. Omitted from this package and not redistributed: 1--641, 650--650, 701--1551. Nothing inside a retained excerpt is omitted or altered: every retained body is delivered exactly as the article's lines are written, so no omission receipt of any kind accompanies this package. Citations: the retained excerpts carry no citation command at all, so no bibliography entry of the article is transcribed and no reference is redistributed. Reference adaptation: none was needed and none was made; every reference inside the delivered unit resolves inside it, so no unresolved reference or citation is produced. Printed slips kept literal: no printed word, symbol, hypothesis or number of the counted statement or of its proof was altered. Layout and class: the article's own class is \texttt{amsart} with packages including geometry, bm, enumerate, amssymb, ulem, color, graphicx, hyperref; the wrapper is the lane's pinned \texttt{amsart} editorial wrapper at 10pt on A4, which restates the theorem-like environments the retained text needs and adds the article's own macro definitions that the retained text needs (\texttt{\textbackslash Im}, \texttt{\textbackslash Re}), with extra wrapper packages (bm, bbm, dsfont, xspace, cancel, mathtools). The delivered numbering is the unit's own. Assets: no retained span contains a figure environment, a graphics-inclusion command, a PDF-page inclusion, a table or a TikZ picture; no image file is copied into this package, the package declares no asset dependency and no image file is redistributed with it. Licence: version arXiv:2105.10711v3 of this article is distributed under the Creative Commons Attribution 4.0 International licence, as recorded in the arXiv licence metadata for this exact version and shown on the abstract page https://arxiv.org/abs/2105.10711v3. A full rights notice is delivered at licenses/arxiv-2105.10711-CC-BY-4.0.txt. Characters: every retained excerpt is pure ASCII, so the delivered engine, XeLaTeX with its default Latin Modern text font, reports no missing character.
\begin{lemma}\label{lm:ReImw}
Let $R_1$ be the portion of the domain $R$ restricted to the first quadrant of the $z$-plane with $w(\infty)=+1$. 
Denote by $R_1^\circ$ the interior of $R_1$. 
For any $z\in R_1^\circ$, we have
\[
\Re(w)>0>\Im(w). 
\]
\end{lemma}

\begin{proof}
In the proof, we always consider $\arg z$ as a single valued function that takes value 
in the interval $(-\pi,\pi]$. 
We first show that the imaginary part of 
\[
w^2=\frac{g_1(z)}{g_2(z)}
   =\frac{(z-\lambda)(z+1)(z-\lambda_1)(z+\lambda_2)}{(z+\lambda)(z-1)(z+\lambda_1)(z-\lambda_2)}
\]
is negative for $z\in R_1^\circ$. 
The map $\mathbb{C}\cup\{\infty\}\ni z\mapsto g_1(z)/g_2(z)\in\mathbb{C}\cup\{\infty\}$ is degree 4 
and the restriction $\mathbb{R}\cup\{\infty\}\ni x\mapsto g_1(x)/g_2(x)\in\mathbb{R}\cup\{\infty\}$ 
is also four to one map. 
Thus $w^2\in\mathbb{R}\cup\{\infty\}$ if and only if $z\in\mathbb{R}\cup\{\infty\}$. 
Hence, by continuity, $\Im(w^2)$ is either positive or negative for all $z\in R_1^\circ$. 
Let $z_0=r+i\epsilon\in R_1^\circ$, where $r$ is sufficiently large positive number and 
$\epsilon$ is sufficiently small positive number such that both $g_1(z_0)$ and $g_2(z_0)$ take 
values in $R_1^\circ$. 
Then there exist $\theta_1,\theta_2\in (0,\pi/2)$ so that 
\[
\arg g_i(z_0) = \theta_i\qquad (i=1,2). 
\]
Then, since 
\begin{align*}
0&<\arg (z_0+\lambda_2)
  <\arg (z_0+\lambda_1)
  <\arg (z_0+1)
  <\arg (z_0+\lambda) \\
 &<\arg (z_0-\lambda)
  <\arg (z_0-1)
  <\arg (z_0-\lambda_1)
  <\arg (z_0-\lambda_2)
  <\frac{\pi}{2}, 
\end{align*}
we have 
\begin{align*}
\theta_1
&=\arg (z_0-\lambda)+\arg (z_0+1)+\arg (z_0-\lambda_1)+(z_0+\lambda_2) \\
&<\arg (z_0-1)+\arg (z_0+\lambda)+\arg (z_0-\lambda_2)+(z_0+\lambda_1)
 =\theta_2,  
\end{align*}
and hence $\arg w^2=\theta_1-\theta_2<0$. 
This implies that $\Im(w^2)<0$. 
Therefore, $w$ satisfies either
\[
\Re(w)<0<\Im(w)
\qquad\mbox{or}\qquad
\Re(w)>0>\Im(w). 
\]
Since $w(\infty)=+1$, we see that $\Re(w)>0>\Im(w)$. 
\end{proof}

\end{document}
